\section{\sys Design}

\subsection{Design Principles}

\sys follows three key design principles to address the limitations of existing GPU stacks:

\paragraph{Principle 1: Decouple stable mechanisms from dynamic programmable policies.}
\sys factors the GPU stack into a small set of stable, privileged mechanisms and a separate, dynamically programmable policy layer. The \sysdriver exposes only those mechanisms that actually manipulate hardware state: page migration, region admission/removal, kernel submission and preemption, and block scheduling, through narrow, versioned interfaces. All higher-level decisions (which region to evict, which kernel to admit next, how to prioritize tenants) live in eBPF policies attached to these hooks. We mandate driver modification because purely user-space or sidecar approaches lack the privileges to interpose on critical hardware paths (e.g., MMU faults, context switching) and cannot enforce global, cross-tenant invariants. This requires one-time driver changes, but once in place it lets us evolve and deploy new policies without touching vendor drivers or GPU runtimes.

\paragraph{Principle 2: Unified CPU--GPU policy interface and control plane.}
All \sys policies, whether running in the OS kernel or inside GPU kernels, are written in the same restricted eBPF IR and use the same abstractions (virtual addresses, time, shared maps). A single control plane loads, verifies, and attaches eBPF programs, configures shared maps, and coordinates their lifecycle. This sacrifices some expressiveness compared to arbitrary CUDA code (e.g., direct kernel manipulations), but eliminates cross-layer impedance mismatch and makes it straightforward to split a policy across host and device.

\paragraph{Principle 3: Safety and efficiency through verified bounded execution and interfaces}
\sys only executes policy code that is expressed in a restricted eBPF subset and accepted by the in-kernel verifier. The verifier enforces bounded loops, stack limits, and memory safety relative to maps and context structs. On the GPU we reuse the same verified bytecode, lower it with a trusted backend, and run it in a sandbox that enforces per-invocation work budgets and restricts memory accesses to \sys-managed regions. This design intentionally trades policy expressiveness for strong guarantees: policy handlers cannot hang the GPU, corrupt driver or application memory, or introduce unbounded latency on critical paths.

\subsection{High-Level Architecture}
\label{sec:arch}

\begin{figure}
\centering
\fbox{\parbox{0.8\columnwidth}{\centering [Figure: \sys Design]}}
\vspace{-.5ex}
\caption{\sys high-level design: unified eBPF runtime across CPU and GPU.}
\label{fig:system-design}
\end{figure}

Figure~\ref{fig:system-design} shows the high-level architecture of \sys, including the CPU--GPU boundary, \sysdriver extensions, and GPU kernel-level runtime integration points. At its core, \sys provides a unified eBPF runtime that manages verification and just-in-time (JIT) compilation in the operating system kernel or the userspace loader.

Policy authors write \sys programs using standard eBPF tooling (e.g., \texttt{clang} with \texttt{libbpf}, \texttt{bpftool}, or \texttt{bpftrace}). Programs are compiled to eBPF bytecode, loaded into the kernel, and verified before execution. \sys places eBPF maps in memory regions visible to both the OS kernel and GPU kernels.

Once verified, the same bytecode can be installed in two ways:

\begin{itemize}[leftmargin=*]
  \item \emph{Host side:} \sys attaches eBPF programs to hooks in \sysdriver at critical policy decision points in the GPU stack (e.g., memory prefetching, migration, and kernel scheduling). These hooks follow a \texttt{struct\_ops}-style callback interface: when the driver reaches a hook, it invokes the attached eBPF handler with an event-specific context structure. The handler returns a compact decision (e.g., an action code or priority) that the driver enforces or using safe kfunc calls into the driver code.
  \item \emph{Device side:} for execution inside GPU kernels, \sys translates the verified eBPF bytecode to device bytecode and injects it into the target GPU context. Lightweight trampolines inserted at kernel entry or other selected sites save minimal live state, call into the translated eBPF snippet, and then restore state to resume the original control flow. Device-side eBPF programs update the same maps as their host-side counterparts.
\end{itemize}

A GPU-side sandbox runtime enforces the safety guarantees described in the Safety Model (\S\ref{sec:prog-model}): each injected program operates in a protected domain that cannot interfere with application kernels, while still enabling adaptive policy decisions on critical GPU execution paths.

Shared eBPF maps provide efficient state coordination between \sysdriver and eBPF policy logic on the host, and between host and GPU policies across the CPU--GPU boundary. This design yields real-time, low-overhead cross-layer observability and control while retaining the safety guarantees of the eBPF framework. We focus on multi-tenant policy control on shared GPUs, not on providing virtual GPU abstractions or hardware partitioning. Mechanisms such as MIG or SR-IOV are orthogonal and can, in principle, be deployed underneath \sys; \sys policies only manipulate well-defined \sys-managed state and cannot corrupt kernel or application memory outside that state.


\subsection{\sys Programming Model}
\label{sec:prog-model}

A \sys policy is a distributed application consisting of eBPF code that runs in
privileged contexts on the CPU and GPU together with a user-space control-plane
component. All pieces program against a unified abstraction: a restricted eBPF IR,
a small set of attach points, and shared eBPF maps and metadata.

\subsubsection{Policy Code (CPU \& GPU)}

\sys policies are expressed as eBPF programs using a unified intermediate representation: a restricted subset of eBPF instructions and helpers. All \sys policies are compiled to standard eBPF bytecode with existing toolchains (clang + libbpf) and loaded into the kernel using unmodified eBPF system calls. The kernel's eBPF verifier checks that each program uses only the \sys-specific kfuncs and hook contexts declared via BTF annotations. We do not implement a separate user-space verifier.

\sys introduces two primary eBPF program types to categorize policy logic:
a driver-level type for GPU policy hooks (encompassing both \texttt{BPF\_PROG\_TYPE\_GPU\_MEM} for memory events and \texttt{BPF\_PROG\_TYPE\_GPU\_SCHED} for scheduling) and a device-level type (\texttt{BPF\_PROG\_TYPE\_GPU\_DEV}) for GPU-kernel instrumentation. Each program type is associated with a specific family of hook contexts that are rigorously defined and BTF-annotated (BPF Type Format) to be understood by the kernel verifier.

Each program defines multiple entry points corresponding to attach points on either the host or GPU device. Policies operate within a strict safety sandbox: they can only read fields explicitly marked as safe within these context types. The verifier enforces that any calls to kernel functions (kfuncs) match their declared prototypes and are permitted for the given program type. Furthermore, policies are restricted from arbitrary writes; they can only modify system state through specialized, safe kfuncs (e.g., updating a calculated prefetch region or setting a scheduling priority), ensuring that policy logic cannot inadvertently corrupt internal driver or device state.

On the host, verified eBPF bytecode executes directly in privileged kernel contexts provided by \sysdriver, allowing policies to respond to critical memory events (such as device memory misses, eviction candidate selection) and kernel scheduling decisions (kernel submission, admission, preemption, completion). On the GPU, the same verified bytecode is translated to device instructions (e.g., PTX or SPIR-V) and executed through lightweight trampolines injected at kernel entry points, selected memory operations, and synchronization or phase boundaries. In both host and GPU contexts, policy handlers receive structured event contexts containing metadata (e.g., virtual addresses, region IDs, kernel IDs, warp identifiers, and resource usage metrics), and return compact policy decisions enforced directly by underlying GPU mechanisms.

\subsubsection{Shared State \& Metadata}

Shared eBPF maps store policy-relevant state including counters, histograms, per-region hotness metrics, per-tenant scheduling priorities, and runtime configuration parameters. This enables device-side probes, such as GPU memory access instrumentation or warp progress signals, to update policy metadata that host-side memory or scheduling hooks later consult. \sys transparently propagates and translates core metadata—such as virtual memory addresses, kernel identifiers, tenant IDs, timestamps, and execution contexts—across CPU and GPU, presenting a logically unified identifier space. This design allows policy developers to implement sophisticated cross-layer logic (e.g., GPU-side access tracking coupled with host-side eviction or scheduling) without manually handling address translation or cross-boundary synchronization complexities.

\subsubsection{Control-Plane Application and Lifecycle}

A \sys policy also includes a dedicated user-space control-plane component responsible for lifecycle management, dynamic reconfiguration, and policy introspection. This component instantiates and manages eBPF maps, and attaches policy sections to defined host and GPU device hook points exposed by \sysdriver and GPU-side runtime. It enables dynamic reconfiguration of deployed policies at runtime by updating parameters (e.g., thresholds, sampling rates, eviction aggressiveness, power caps) stored in shared maps, and facilitates atomic swapping or updating of attached policy logic without restarting applications or kernels. Furthermore, the control-plane application continuously monitors policy operation by reading metrics and state from maps—such as memory hotness counters, per-tenant scheduling metrics, and queueing latencies—and exports this information to external observability frameworks or adaptation controllers. Thus, from a developer's perspective, a complete \sys policy application comprises policy logic for both CPU and GPU execution contexts along with a complementary control-plane component orchestrating policy behavior and evolution over time.

\paragraph{Safety Model and TCB.}
The core challenge in programmable GPU policy is to run operator-supplied logic in privileged, latency-critical GPU driver and device paths without process-like isolation: even simple bugs can corrupt mappings, stall the device, or inflate tail latency. \sys runs operator-supplied policy code in these privileged paths. We assume the entity deploying \sys policies (e.g., a cluster operator) is trusted but fallible: policies may contain logic bugs or misconfigurations but are not malicious. Malicious kernel modules, vendor drivers, firmware, or hardware are out of scope.

Under this model, the \sys TCB consists of: (1) the in-kernel eBPF verifier, (2) the \sys kernel module, (3) the device-code backend, and (4) the GPU driver/firmware. Policy authors only reason about eBPF and the \sys API; they never write PTX/SASS directly. All policies are expressed in a restricted eBPF subset and must pass the in-kernel verifier, which enforces bounded loops, stack and map safety, and restricts helpers and kfuncs to the \sys-specific set.

\sys reuses the verified bytecode on the GPU: a trusted backend lowers it to device code, and a SIMT-aware runtime executes one lane per warp in a sandboxed context. We also use conservative work budgets for execution instructions and memory access based on verifier bounds. The sandbox derives an instruction budget from verifier bounds, restricts memory accesses to \sys-managed maps and per-warp scratch (with explicit bounds checks), and discards results if the budget is exceeded. We do not run a separate device-side verifier; instead, the backend is part of the TCB and cannot grant capabilities beyond those allowed by the verifier. This gives operators something prior GPU tools lack: a driver-anchored policy plane where the same verified program can run in the GPU driver and inside GPU kernels, with verifier-backed memory safety and bounded per-hook work.



\paragraph{Optimizing for GPU's SIMT Execution.}
Naively executing scalar control logic across all GPU threads introduces prohibitive divergence and overhead. To overcome this, \sys aggregates scalar policy logic execution at warp granularity, executing decisions on a designated lane while synchronizing minimal input from other lanes. Frequent global state updates are optimized through warp-level and SM-level local aggregation techniques, dramatically reducing contention and costly global memory operations. Context metadata passed to policy handlers is minimized through static analysis, passing only essential fields to reduce memory traffic and latency. Adaptive sampling mechanisms allow policies to dynamically tune instrumentation frequency, achieving low overhead across diverse workloads. Collectively, these optimizations enable efficient, scalable, and minimally invasive GPU-side policy execution suitable for production deployment.

\subsection{Policy Interface for Memory}
\label{sec:design-mem}

Modern GPU memory subsystems move data across device memory, host memory, and
intermediate tiers using fixed, internally implemented heuristics, often tied to
page faults or simple recency-based policies. These mechanisms neither expose their
decisions nor exploit workload-specific structure, and they interact poorly with
user-space caches that increasingly manage model weights, KV-cache segments, and
expert parameters. \sys instead treats memory placement as a programmable region
cache: it exposes a small set of region-level events for \emph{admission},
\emph{access}, and \emph{removal}, together with a prefetch channel, and lets host
and device cooperate on a unified policy.

On the host, these events are realized as a small \texttt{struct\_ops}-style policy
table with typed contexts:

\begin{minted}[highlightlines=false,fontsize=\small]{c}
// Host-side policy callbacks for GPU memory management.
struct gdrv_mem_ops {
  int (*region_add)(struct memsys_region_add_ctx *ctx);
  int (*region_access)(struct memsys_region_access_ctx *ctx);
  int (*region_remove)(struct memsys_region_remove_ctx *ctx);
  int (*prefetch)(struct memsys_prefetch_ctx *ctx);
  char name[MEMSYS_NAME_LEN];
};

// device-side memory callback interface
struct gdev_mem_ops {
    void (*memtrace)(...);
    ...
};
\end{minted}

\paragraph{Host-side region events and removal.}
On the host, \sys models memory as a set of logical regions (for example, weight
shards, KV-cache segments, or embedding tables) and exposes attach points for three
canonical cache events:

\begin{itemize}[leftmargin=*]
  \item \emph{Add}, when a region is created or becomes eligible for device placement.
        The context includes a region identifier, size, tenant identifier, and
        optional semantic tags (for example, read-only, streaming, or shared).
  \item \emph{Access}, when the memory subsystem observes use of a region and must
        decide whether to treat it as a cache hit, trigger migration, or bypass
        the cache altogether. Access events may be driven by hardware misses,
        by explicit runtime calls, or by higher-level cache managers built on top
        of \sys.
  \item \emph{Remove}, when a region is about to leave the device-resident working
        set. The context includes a cause flag and a candidate set. For the
        \emph{lifetime} case, an upper layer has logically freed the region and
        the candidate set is a singleton; the policy can discard any associated
        state and reclaim capacity. For the \emph{pressure} case, the memory
        subsystem has determined that some regions must be evicted under capacity
        constraints and supplies a set of candidates; the policy ranks these
        candidates and selects preferred victims to remove.
\end{itemize}

In addition, \sys provides a \emph{prefetch hook} that runs at safe points where
data can be moved proactively, for example between kernel launches or at explicit
scheduler ticks. The prefetch context includes a snapshot of current region state,
available capacity, and any device-generated hints, and the policy returns a ranked
list of regions to promote to device memory. Together, these events realize a
general region-cache interface: policies implemented as eBPF handlers maintain their
own state in shared maps (for example, LRU lists, LFU counters, or tenant budgets),
and respond to Add/Access/Remove and prefetch callbacks with compact decisions
such as \emph{admit}, \emph{bypass}, \emph{pin}, \emph{evict}, or ordered lists
of regions to migrate. The underlying memory subsystem enforces these decisions
and maintains correct mappings, but delegates replacement and prefetch policy to
\sys.

\paragraph{Unified device-side memory events.}
On the GPU, \sys exposes a single device-side memory event interface that kernels
use both to report observed accesses and to express future-use hints. Device-side
handlers receive a compact context struct (not shown) with fields such as a logical
region identifier, kernel, warp and SM identifiers, an event type (\emph{ACCESS}
or \emph{HINT}), and optional hint flags.

At selected memory operation sites, \sys injects lightweight hooks that operate at
warp granularity and invoke a device-side \sys handler with an \emph{ACCESS} event:
the handler updates shared maps with per-region statistics such as access counts,
approximate working-set membership, and simple locality signatures. To bound
overhead, \sys supports sampling (for example, only one in every $N$ accesses
generates an event) and per-warp aggregation before these statistics are written
to shared maps.

The same device-side interface also supports \emph{HINT} events, which kernels
invoke explicitly when they have domain knowledge about future region usage. For
example, a mixture-of-experts kernel can compute the identifiers of experts it
will access for the next batch and call a \sys helper with HINT events that mark
the corresponding regions as “will-use soon” with an optional priority or deadline;
a KV-cache kernel can emit HINT events that mark segments as cold and suitable for
demotion. Both ACCESS and HINT events are recorded in shared maps or small control
queues that host-side \sys handlers consult at Access, Remove, and prefetch hooks.
Device-side code never manipulates page tables or performs migration directly; it
implements the sensing and hinting side of a closed-loop policy, while host-side
\sys handlers implement the placement, prefetch, and eviction decisions over the
GPU memory hierarchy.


\subsection{Policy Interface for Scheduling}
\label{sec:design-sched}

GPU schedulers today typically implement fixed, coarse-grained policies such as FIFO
queues with limited priority control, which are insufficient in multi-tenant settings
with mixed latency-sensitive and batch workloads or tight power constraints. In such
environments, operators need policies that can shape both \emph{which kernels reach
the GPU} and \emph{how work is consumed inside those kernels} at the granularity of
thread blocks or tiles. \sys exposes scheduling interfaces on both host and device
so that policies can implement fine-grained priority, deadline-aware, and power-aware
scheduling while respecting the GPU SIMT execution model.

On the host, these lifecycle events are exported as a small struct-ops table of
scheduling callbacks:

\begin{minted}[highlightlines=false,fontsize=\small]{c}
// Host-side policy callbacks for GPU kernel scheduling.
struct gdrv_sched_ops {
  // Called when a kernel is enqueued by a runtime.
  // Return 0 to accept, <0 to reject or defer.
  int  (*submit)(struct gpusched_submit_ctx *ctx);
  // Called when the scheduler is about to dispatch a ready kernel.
  // Returns the ID of the kernel to admit, or -1 to defer.
  int  (*admit)(struct gpusched_admit_ctx *ctx);
  // Called at preemption opportunities. Returns an action code:
  // e.g., GDRV_SCHED_CONTINUE or GDRV_SCHED_PREEMPT.
  int  (*preempt)(struct gpusched_preempt_ctx *ctx);
  // Notification that a kernel has completed. Must not fail.
  void (*complete)(struct gpusched_complete_ctx *ctx);
  char name[GPUSCHED_NAME_LEN];
};
// Device-side block scheduling callbacks.
struct gdev_sched_ops {
  // Called when a worker block starts a new work unit.
  void (*entry)(struct gdev_block_ctx *ctx);
  // Called when a worker block finishes a work unit.
  void (*exit)(struct gdev_block_ctx *ctx);
  // Called when a worker block is idle; returns whether
  // it should try to steal work and from which tenant/queue.
  bool (*should_try_steal)(struct gdev_block_ctx *ctx);
};
\end{minted}

\paragraph{Host-side scheduling hooks.}
On the host, \sys extends the GPU kernel lifecycle with attach points that allow
policies to control admission, ordering, and preemption of kernels:

\begin{itemize}[leftmargin=*]
  \item \texttt{submit} runs when a framework enqueues a kernel; the context includes kernel metadata (grid and block size), tenant or stream identifiers, and priority hints. Policies may reject kernels (e.g., enforcing admission control) or adjust their priority fields before returning.
  \item \texttt{admit} runs when the driver is about to dispatch a kernel to a hardware queue; it can choose among ready kernels or defer admission by returning \texttt{-1}, effectively implementing arbitrary queueing disciplines over driver queues.
  \item \texttt{preempt} runs at cooperative yield points or at the end of a time slice; it returns a simple action code (e.g., \texttt{CONTINUE} vs \texttt{PREEMPT}), which the driver enforces by marking the kernel preemptible and enqueueing higher-priority work.
  \item \texttt{complete} is a notification that cannot fail; policies use it to update per-tenant statistics, virtual time, or longer-term admission control state.
\end{itemize}

Attached eBPF schedulers read these contexts and shared maps (for example, per-tenant
latency and throughput statistics, outstanding work across queues) and return compact
decisions such as which kernel to run next, whether to throttle a tenant, or whether
to allow preemption. The host-side interface thus controls which kernels occupy GPU
contexts and how they are multiplexed in multi-tenant scenarios.

\paragraph{Device-side thread-block scheduler and signals.}
Inside the GPU, \sys exposes a complementary interface that operates at the level of
thread blocks or tiles, aligning with the SIMT execution model. Each kernel that uses \sys's internal scheduler is modeled as a pool of \emph{work units} (e.g., tiles or logical block indices) consumed by \emph{worker blocks}. Device-side policies observe and steer this internal scheduler via a \texttt{gdev\_block\_ctx} that includes the kernel identifier, tenant identifier, and a logical work-unit or block identifier.

\begin{itemize}[leftmargin=*]
  \item \texttt{entry} is invoked (one lane per warp) when a worker block starts a new work unit. Policies can record progress or adjust per-tenant counters in BPF maps.
  \item \texttt{exit} runs when a work unit completes and can update per-tenant service metrics or mark a block as safe to preempt.
  \item \texttt{should\_try\_steal} is called when a worker becomes idle; it reads shared maps and returns whether to pull work from another queue (e.g., a high-priority tenant) or park the block.
\end{itemize}

In addition, \sys injects lightweight probes at kernel entry and at selected synchronization or phase boundaries to record per-warp or per-block progress (for example, iterations completed in a loop, entry into a barrier-heavy phase) and simple behavioral signals such as whether a block is currently memory-bound or compute-bound. These probes write into shared maps and can mark when a block or phase is safe to preempt.

All device-side handlers are sandboxed and must complete within the verifier-derived budget. They cannot directly manipulate hardware queues; instead, they communicate hints, progress information, and per-block decisions that host policies or the \sys runtime consume at the next host-side scheduling hook or device scheduling point. Together, the host-side and device-side interfaces allow \sys policies to coordinate decisions across the full scheduling pipeline: from kernel admission and preemption on the host to fine-grained block-level work distribution inside GPU kernels.

\paragraph{Example: priority-aware block scheduler.}
As a concrete illustration, consider a policy that gives low latency to an interactive tenant while preserving throughput for batch jobs. The host keeps a per-tenant record of outstanding work and virtual time in a BPF map. \texttt{submit} tags kernels with a tenant ID and initial priority. \texttt{admit} prefers kernels from the interactive tenant when its queue is non-empty and boosts its priority when its p99 latency (tracked in another map) exceeds a target. On the device, each worker block's \texttt{entry} handler increments a per-tenant ``active blocks'' counter, and \texttt{exit} decrements it and records the work-unit completion time. When \texttt{should\_try\_steal} runs on an idle block, it uses these maps to decide whether to steal a work unit from the interactive tenant or from the batch pool. In \S\ref{sec:eval} we show that this policy reduces p99 latency by 2--5$\times$ with 8--12\% throughput loss for batch work.
